\chapter{Starting from Zero}
\label{ch:zero}

This course has no prerequisites. Most of its students have never written a
program, and none is expected to know how a computer is built. This chapter is
for them. It explains, in plain words, the few ideas and the few dozen words
that the first lectures use from the very first day; everything here is taught
again, properly, in its own chapter. Read it before Lecture~01, and come back
to it whenever a word on a slide is unfamiliar: the small dictionary at the end
says where each one is explained in full.

\section*{What a computer does}

A computer has three kinds of parts. The \term{processor} does the work. The
\term{memory} keeps the program and the data while the computer is on. The
\term{input and output} devices connect it to the world: a keyboard, a
joystick, a screen, a speaker, a network.

The processor does one thing, very fast: it executes \term{instructions}.
An instruction is a small, exact order, such as ``add these two numbers'',
``copy this value from memory'', or ``if this number is zero, continue from
over there''. Each instruction does very little; a processor executes hundreds
of millions of them every second, and everything a computer does, from a game
to a video call, is a long sequence of such small steps. A \term{program} is
the list of instructions that produces a given behaviour. Chapter~\ref{ch:history}
tells how this idea, a program kept in memory like any other data, was born in
1945, and why it still describes the computer on your desk.

\section*{Bits, bytes and the numbers on the page}

Inside a computer everything, numbers, letters, colours, instructions, is
stored as \term{bits}: digits that can only be 0 or 1. Eight bits form a
\term{byte}; thirty-two bits, four bytes, form a \term{word}, the size of the
numbers that the processor of this course handles in one step.

A number written with bits is in \term{binary}, base~2: each position is worth
twice the one on its right, so $1010_2 = 8 + 2 = 10$. Binary is exact but
long, so programmers group the bits four at a time and write each group as one
\term{hexadecimal} digit, base~16: the digits 0 to 9 and then A to F for ten
to fifteen. The prefix \texttt{0x} marks a hexadecimal number.

\begin{center}\small
\begin{tabular}{rll}
\toprule
\textbf{Decimal} & \textbf{Binary} & \textbf{Hexadecimal}\\
\midrule
0 & \texttt{0000} & \hex{0}\\
1 & \texttt{0001} & \hex{1}\\
9 & \texttt{1001} & \hex{9}\\
10 & \texttt{1010} & \hex{A}\\
15 & \texttt{1111} & \hex{F}\\
16 & \texttt{1 0000} & \hex{10}\\
255 & \texttt{1111 1111} & \hex{FF}\\
\bottomrule
\end{tabular}
\end{center}

\noindent Because one hexadecimal digit is exactly four bits, the colour green
of \prg, \hex{07E0}, is \texttt{0000 0111 1110 0000} in binary. For the first
weeks it is enough to recognise these notations; Chapter~\ref{ch:information}
explains them, with conversions in both directions.

\section*{From text to machine code}

The processor reads instructions as numbers. This numeric form of a program is
called \term{machine code}. Nobody writes it by hand: people write
\term{source code}, a text file in a \term{programming language}, and a
translator turns it into machine code. This course uses \term{assembly}, the
language closest to the machine, in which one line becomes one instruction; the
translator for assembly is the \term{assembler}. For C, a language that the
\prg examples also use, the translator is a \term{compiler}. The set of these
translators and of their helpers is the \term{toolchain}, which Lecture~01
installs. Here is one instruction in its two forms:

\begin{center}\small
\begin{tabular}{ll}
\toprule
\textbf{You write (source code)} & \textbf{The processor reads (machine code)}\\
\midrule
\texttt{addi a0, a0, 1} & \hex{00150513}\\
\bottomrule
\end{tabular}
\end{center}

\noindent On \prg the result of the translation is a \term{cartridge}: one
small file, ending in \fname{.prg32}, that holds the machine code of a game, its
data and a short header, like the cartridges of the consoles of the 1980s.

\section*{Who runs your program}

A game never runs on an empty machine. On the \prg kit, another program, the
\term{firmware}, starts when the board is switched on, stays in memory and owns
the screen, the buttons and the sound. It is kept in \term{flash} memory, which,
like a USB stick, keeps its contents when the power is off. The part of the
firmware that a game talks to is the \term{runtime}: it calls the game once for
every \term{frame}, every new picture on the screen, and lends it ready-made
functions to draw and to play sounds.

Any machine that offers this runtime is a \term{host}. There are three: the
board; QEMU, an \term{emulator}, a program that imitates a processor on a
laptop; and \prgqt, an application whose \term{interpreter} reads each
instruction of the cartridge and carries it out in software. The same
cartridge runs on all three. \prgqt also contains the \term{debugger} of the
course: a tool that pauses a program, runs it one instruction at a time, and
shows what is inside the machine at every step.

\section*{The terminal, in ten minutes}

Installing and building use the \term{terminal}: a window in which commands are
typed instead of clicked. The terminal shows a \term{prompt}, a short line
ending in \texttt{\$} or \texttt{>}, and waits. One types a command and presses
Enter; the computer runs it and prints its answer below. At every moment the
terminal is ``in'' a folder, the \term{current folder}, and commands act on
the files there. A \term{path} names a file by the folders that lead to it,
such as \fname{\textasciitilde/PRG32/}\allowbreak\fname{examples/pong.S}, where
\fname{\textasciitilde} is your home folder.

\begin{center}\small
\begin{tabular}{lll}
\toprule
\textbf{To} & \textbf{macOS and Linux} & \textbf{Windows PowerShell}\\
\midrule
see which folder you are in & \texttt{pwd} & \texttt{pwd}\\
list the files in that folder & \texttt{ls} & \texttt{dir}\\
enter the folder PRG32 & \texttt{cd PRG32} & \texttt{cd PRG32}\\
go back up one folder & \texttt{cd ..} & \texttt{cd ..}\\
go to your home folder & \texttt{cd \textasciitilde} & \texttt{cd \textasciitilde}\\
stop the program that is running & Ctrl+C & Ctrl+C\\
\bottomrule
\end{tabular}
\end{center}

\noindent The Tab key completes a name that you have started to type, and the
up arrow recalls the previous command. A \term{repository} is a folder of
source code whose history is kept by a program called \texttt{git};
\texttt{git clone} copies one from the Internet onto your laptop.

A command that works often prints little or nothing. A command that fails
prints a message, and the message almost always names the cause: read all of
it, from the first line that says \emph{error}, before changing anything.
Four messages cover most problems of the first week. \emph{command not found}:
the program is not installed, or not active in this terminal. \emph{No such file
or directory}: wrong folder or wrong name, which \texttt{pwd} and \texttt{ls}
reveal. \emph{Permission denied}: the command needs rights that you do not
have here. \emph{Connection refused}: the program you are talking to, such as
\prgqt, is not running. When you ask for help, copy the command and the whole
message as text.

\section*{Reading one line of assembly}

A line of assembly has at most four parts:

\begin{asm}
first_update:   addi t1, t1, 1      # one more tick
\end{asm}

\noindent \code{first\_update:} is a \term{label}, a name for this place in the
program, ending in a colon; a label often marks the start of a
\term{function}, a piece of program with a name that does one job.
\ins{addi} is the \term{instruction}, here ``add a constant''. \code{t1, t1, 1}
are the \term{operands}: the destination comes first, then the sources, so the
line reads ``put \reg{t1} plus 1 into \reg{t1}''. Everything after \code{\#} is
a \term{comment}, written for people and ignored by the assembler. A line that
begins with a dot, such as \code{.section .text}, is a \term{directive}: an
order to the assembler, not to the processor.

\section*{Registers, memory, addresses}

The processor computes only on its \term{registers}: thirty-two small storage
cells inside it, each holding one word. Memory is much larger: a long row of
cells, one byte each, numbered from zero; the number of a cell is its
\term{address}. To work on a value kept in memory, a program first
\term{loads} it into a register (\ins{lw}), computes on it, and then
\term{stores} it back (\ins{sw}). The difference between an address, where a
value lives, and the value itself is the single most important idea of the
first weeks; Chapter~\ref{ch:firstcontact} shows it with two arrows.

To use a function, a program \term{calls} it: \ins{call} jumps to the first
instruction of the function and records where to come back in the register
\reg{ra}; \ins{ret} returns there. The values that the function needs, its
\term{arguments}, travel in the registers \reg{a0} to \reg{a7}.

\section*{A small dictionary for the first weeks}

\begin{description}\small
\item[ABI] the agreement on how a program and the runtime call each other; Chapters~\ref{ch:setup} and~\ref{ch:processor}.
\item[Bit, byte, word] one binary digit; eight bits; thirty-two bits; Chapter~\ref{ch:information}.
\item[Breakpoint] a mark on an instruction where the debugger stops the program; Chapter~\ref{ch:setup}.
\item[Clock] the signal that paces the processor; 160\,MHz means 160 million steps per second; Chapter~\ref{ch:sequential}.
\item[Frame] one picture on the screen; a game draws thirty or more each second; Chapter~\ref{ch:firstcontact}.
\item[Instruction set (ISA)] the list of instructions a processor understands; Chapters~\ref{ch:history} and~\ref{ch:instructions}.
\item[Pixel] one coloured dot of the screen; Chapter~\ref{ch:information}.
\item[Stack] an area of memory where functions keep values for a while; Chapters~\ref{ch:firstcontact} and~\ref{ch:processor}.
\item[Two's complement] the way computers store negative integers; Chapter~\ref{ch:integers}.
\end{description}

\noindent Appendix~\ref{app:glossary} lists every term of the course, in
English and in Italian.

\section*{How to study when you start from zero}

Three habits matter more than any previous knowledge. \emph{Predict before you
run}: before executing a program, write down what you expect to see; a wrong
prediction is the most useful thing that can happen, because it shows exactly
what to learn next. \emph{Read the whole message}: errors are information, not
failure, and almost every one names its cause. \emph{Ask early}: a question
asked in the first weeks costs minutes, the same question in December costs
days. The questions at the end of each chapter take a few minutes a week; they
are the best preparation for the examination, and the best way to find out,
early, what is not yet clear.
